% --- 1. Essential Packages ---
\usepackage[utf8]{inputenc}
\usepackage{amsmath, amssymb, amsthm} % The Holy Trinity of math typesetting
\usepackage{mathtools}                % Provides \coloneqq for definition (:=)
\usepackage{enumitem}                 % Better lists (e.g., (a), (b), (c))
\usepackage{geometry}                 % Easily adjust page margins
\geometry{margin=1in}

% --- 2. Number Sets ---
\newcommand{\R}{\mathbb{R}}
\newcommand{\N}{\mathbb{N}}
\newcommand{\Z}{\mathbb{Z}}
\newcommand{\Q}{\mathbb{Q}}
\newcommand{\C}{\mathbb{C}}

% --- 3. Real Analysis Variables ---
% \varepsilon looks much better and is the standard over \epsilon
\newcommand{\eps}{\varepsilon}
\newcommand{\del}{\delta}

% --- 4. Auto-sizing Brackets ---
% These automatically resize to fit whatever fraction or term is inside them
\newcommand{\set}[1]{\left\{ #1 \right\}}  % \set{x \in \R \mid x > 0}
\newcommand{\abs}[1]{\left| #1 \right|}    % \abs{x - y} < \eps
\newcommand{\norm}[1]{\left\| #1 \right\|} % \norm{f}

% --- 5. Math Operators ---
% Proper spacing for functions (e.g., \sup, \inf are built-in, but these are not)
\DeclareMathOperator{\cl}{cl}       % Closure
\DeclareMathOperator{\interior}{int} % Interior (note: \int is already integral!)
\DeclareMathOperator{\bd}{bd}       % Boundary
\DeclareMathOperator{\diam}{diam}   % Diameter
\DeclareMathOperator{\dom}{dom}     % Domain

% --- 6. Document Formatting ---
% Stop indenting new paragraphs and add a little space between them
\setlength{\parindent}{0pt}
\setlength{\parskip}{0.8em}

% --- 7. Environments (Notes & Theorems) ---
\theoremstyle{definition} % Bold title, normal text body
\newtheorem{definition}{Definition}[section]
\newtheorem{example}{Example}[section]
\newtheorem{problem}{Problem}

\theoremstyle{plain} % Bold title, italicized text body (for formal logic)
\newtheorem{theorem}{Theorem}[section]
\newtheorem{lemma}[theorem]{Lemma}
\newtheorem{corollary}[theorem]{Corollary}
\newtheorem{proposition}[theorem]{Proposition}

% --- 8. Problem Set Specifics ---
% A custom "Solution" environment that puts a nice QED box at the end
\newenvironment{solution}
  {\begin{proof}[Solution]}
  {\end{proof}}
